\subsection{\texorpdfstring{向量空间 \(R^{n}\)}{向量空间 R 的 n 次幂}}

由于 \(\mathbf{R}\) 是一个域，我们可以在 \(\mathbf{R}^n\) 上定义一个域 \(\mathbf{R}\) 上的向量空间结构，乘积 \(tx\) ——标量 \(t \in \mathbf{R}\) 与点（或向量）\(\boldsymbol{x} = (x_i)\) （\(\mathbf{R}^n\) 的）之积——是点 \((tx_i)\)。注意位似 \((t, x) \to tx\) 在 \(\mathbf{R} \times \mathbf{R}^n\) 上是连续的。若 \(\boldsymbol{e}_i\) 表示 \(\mathbf{R}^n\) 中除指标 \(i\) 的坐标等于 \(1\) 外其余坐标全为零的向量，则诸 \(\boldsymbol{e}_i\) 构成向量空间 \(\mathbf{R}^n\) 的一个基，称为该空间的典范基。每个向量 \(\boldsymbol{x} = (x_i) \in \mathbf{R}^n\) 都可写成 \(\boldsymbol{x} = \sum_{i=1}^{n} x_i \boldsymbol{e}_i\) ，而且关系 \(\sum_{i=1}^{n} t_i \boldsymbol{e}_i = 0\) 蕴含 \(t_i = 0\) ，对 \(1 \leqslant i \leqslant n\) 。

因此，在代数的意义上（《代数》第 II 章，§ 7，第 2 小节），向量空间 \(\mathbf{R}^n\) 是 \(n\) 维的（在域 \(\mathbf{R}\) 上）；由此得到它的名称：\(n\) 维实数空间。

设 \(f\) 是向量空间 \(\mathbf{R}^n\) 到向量空间 \(\mathbf{R}^m\) 中的一个仿射映射（ \(m\) 与 \(n\) 为 \(>0\) 的整数）。若令 \(g(\boldsymbol{x}) = f(\boldsymbol{x}) - f(0)\) ，则 \(g\) 是 \(\mathbf{R}^n\) 到 \(\mathbf{R}^m\) 中的线性映射。设 \(a_{ij}\) （ \(1 \leqslant j \leqslant m\) ）是 \(g(e_i)\) 在 \(\mathbf{R}^m\) 中的坐标，\(b_j\) （ \(1 \leqslant j \leqslant m\) ）是 \(f(0)\) 的坐标；若 \(x_i\) （ \(1 \leqslant i \leqslant n\) ）是第 \(i\) 个坐标（\(\boldsymbol{x} \in \mathbf{R}^n\) 的），而 \(y_j\) 是第 \(j\) 个坐标（\(y = f(\boldsymbol{x})\) 的），则我们有

\[
y _ {j} = \sum_ {i = 1} ^ {n} a _ {i j} x _ {i} + b _ {j} \quad (1 \leqslant j \leqslant m).
\]

由于 \(x_1, x_2, \ldots, x_n\) 的每个线性多项式在 \(\mathbf{R}^n\) 上是一致连续的，由此推出 \(\mathbf{R}^n\) 到 \(\mathbf{R}^m\) 中的每个仿射映射在 \(\mathbf{R}^n\) 上是一致连续的（第 II 章，§ 2，第 6 小节，命题 7）。

特别地，我们知道 \(R^{n}\) 到其自身的每个仿射映射都是双射，而且它的逆仍是仿射映射；因此 \(R^{n}\) 到其自身的每个仿射映射是同胚（并且是 \(R^{n}\) 的一致结构的一个自同构）。

设 \((\pmb{a}_i)_{1\leqslant i\leqslant n}\) 是 \(n\) 个向量（\(\mathbf{R}^n\) 中的）组成的一个自由系{[}换言之，向量空间 \(\mathbf{R}^n\) 的一个基{]}；若 \(\pmb{b}\) 是 \(\mathbf{R}^n\) 中任意一点，则集 \(\mathbf{P}\) 由这样的点 \(\pmb{x} = \pmb{b} + \sum_{i=1}^{n} u_i a_i\) 组成：\(-1 \leqslant u_i \leqslant 1\) （对 \(1 \leqslant i \leqslant n\)）；它是 \(\pmb{b}\) 的一个紧邻域；因为存在一个双射仿射映射 \(f\) （\(\mathbf{R}^n\) 到其自身的），使得 \(f(\pmb{b}) = 0\) 且 \(f(\pmb{b} + a_i) = e_i\) （对 \(1 \leqslant i \leqslant n\)）；而 \(f(\mathbf{P})\) 是 \(n\) 个区间 \([-1, +1]\) （在 \(n\) 个因子空间中）的积所成的立方体。\(\mathbf{P}\) 称为以 \(\pmb{b}\) 为中心

、以 \(a_{i}\) 为基向量的闭超平行体。P 的内部由这样的点 \(b + \sum_{i=1}^{n} u_{i} a_{i}\) 组成：\(-1 < u_{i} < 1\) ，对 \(1 \leqslant i \leqslant n\) ；它称为以 b 为中心、以 \(a_{i}\) 为基向量的开超平行体。

